\subsection{Normalización de la complejidad teórica}\label{sec:normalizacion}

La notación Big-$O$ describe crecimiento asintótico, pero no da tiempos en segundos: la expresión
de $T_p$ es adimensional y no puede superponerse directamente sobre las mediciones experimentales.
\emph{Normalizar} significa introducir las constantes de proporcionalidad que faltan para convertir
esa expresión en segundos, calibrándolas con los propios datos medidos. Con ello puede graficarse
$T_{\text{theory}}(p)$ junto a $T_{\text{exp}}(p)$ y verificar que ambas siguen la misma forma.

\subsubsection{Modelo con constantes físicas}

Se reescribe el tiempo total introduciendo tres constantes: $C$ (tiempo por unidad de trabajo de
cómputo), $\alpha$ (latencia por mensaje) y $\beta$ (tiempo por byte transmitido):
\[
T(p)
=
C\left[
\frac{n_{tr}n_{te}}{p}(d+\log_2 K) + n_{te}K\log_2 p
\right]
+
p\alpha
+
4\beta n_{tr}d
+
(\alpha + 4\beta n_{te}d)\log_2 p
+
(\alpha + 8\beta n_{te}K)\log_2 p.
\]

\subsubsection{Instanciación con los parámetros del experimento}

Para el caso de escalabilidad fuerte con $n_{tr}=100000$, $n_{te}=5000$, $d=64$ y $K=3$
(dataset \texttt{digits}, imágenes $8\times8$), se evalúan los coeficientes numéricos
(con $d+\log_2 K = 65.585$):
\[
n_{tr}n_{te}(d+\log_2 K) = 3.279\times10^{10},
\qquad
n_{te}K = 1.5\times10^4,
\]
\[
4n_{tr}d = 2.56\times10^7,
\qquad
4n_{te}d = 1.28\times10^6,
\qquad
8n_{te}K = 1.2\times10^5.
\]
Sustituyendo y agrupando los dos términos de comunicación en $\log_2 p$:
\[
\boxed{
T(p)
=
C\left(
\frac{3.279\times10^{10}}{p}
+
1.5\times10^4\log_2 p
\right)
+
2.56\times10^7\,\beta
+
\alpha\,p
+
\left(2\alpha + 1.40\times10^6\,\beta\right)\log_2 p
}.
\]
Las únicas incógnitas son la variable $p$ y las constantes $C$, $\alpha$ y $\beta$.

\subsubsection{Forma canónica y recuperación de constantes}

La expresión anterior se ajusta a un modelo de cuatro formas funcionales:
\[
T(p) = a_0 + \frac{a_1}{p} + a_2\log_2 p + a_3\,p,
\]
donde cada coeficiente concentra un mecanismo físico distinto:
\[
a_0 = 2.56\times10^7\,\beta
\quad(\text{reparto del entrenamiento}),
\qquad
a_1 = 3.279\times10^{10}\,C
\quad(\text{cómputo}),
\]
\[
a_2 = 1.5\times10^4\,C + 2\alpha + 1.40\times10^6\,\beta
\quad(\text{bcast + reduce}),
\qquad
a_3 = \alpha
\quad(\text{latencia } p\alpha).
\]
Los coeficientes $a_0,\dots,a_3$ se estiman por mínimos cuadrados a partir de los tiempos medidos
en $p = 1,2,4,8,16,32$. De ellos se recuperan las constantes físicas:
\[
\alpha = a_3,
\qquad
\beta = \frac{a_0}{2.56\times10^7},
\qquad
C = \frac{a_1}{3.279\times10^{10}}.
\]

\subsubsection{Validación}

Como se ajustan cuatro coeficientes pero solo hay tres constantes físicas, el coeficiente $a_2$
queda como \emph{ecuación de validación}: debe cumplirse, dentro del error de ajuste,
\[
a_2 \approx 1.5\times10^4\,C + 2\alpha + 1.40\times10^6\,\beta.
\]
Si la relación se satisface y la curva $T(p)$ ajustada reproduce los puntos experimentales, se
concluye que el modelo de complejidad captura correctamente el comportamiento observado. En
particular, el término $a_3\,p$ explica por qué el tiempo deja de mejorar e incluso aumenta para
$p$ grande, mientras que $a_1/p$ domina en la región de escalado casi lineal.

\subsubsection{Ajuste para la versión v2}\label{sec:norm-v2}

Se aplica el procedimiento a la versión v2 (descomposición por entrenamiento con árboles
explícitos) en el caso de escalabilidad fuerte con $n_{tr}=100000$ y $n_{te}=5000$. Los tiempos
totales medidos son:

\begin{table}[H]
\centering
\begin{tabular}{rrrrrrr}
\toprule
$p$ & 1 & 2 & 4 & 8 & 16 & 32 \\
\midrule
$T_p^{exp}$ (s) & 11.119 & 5.835 & 3.028 & 2.088 & 1.944 & 2.022 \\
\bottomrule
\end{tabular}
\caption{Tiempos totales de v2 usados en el ajuste ($n_{tr}=100000$, $n_{te}=5000$, $d=64$, $K=3$).}
\end{table}

Como el término constante de reparto del entrenamiento ($a_0$) resulta indistinguible con solo seis
puntos (el ajuste sin restricción lo lleva a un valor negativo, no físico), se emplea el modelo de
tres formas
\[
T(p) = \frac{a_1}{p} + a_2\log_2 p + a_3\,p,
\]
cuyas constantes quedan exactamente determinadas. El ajuste por mínimos cuadrados da:
\[
a_1 = 11.087,
\qquad
a_2 = 0.1608,
\qquad
a_3 = 0.0289,
\qquad
R^2 = 0.9991.
\]
Recuperando las constantes físicas mediante $C = a_1/(3.279\times10^{10})$, $\alpha = a_3$ y
$\beta = (a_2 - 1.5\times10^4\,C - 2\alpha)/(1.40\times10^6)$:
\[
C \approx 3.38\times10^{-10}\ \text{s/unidad},
\qquad
\alpha \approx 28.9\ \text{ms},
\qquad
\beta \approx 7.36\times10^{-8}\ \text{s/byte}.
\]
Las tres son positivas y del orden esperado (son constantes \emph{efectivas} que absorben el
\emph{overhead} de \texttt{mpi4py} y las barreras de sincronización).

La descomposición del tiempo ajustado muestra el cruce entre cómputo y comunicación: en $p=1$ el
cómputo ($a_1/p$) es el $99.7\%$ del tiempo, pero su peso cae hasta $\sim38\%$ en $p=16$ y
$\sim17\%$ en $p=32$, punto en que la comunicación ($a_2\log_2 p + a_3\,p$) ya domina. Esto valida
la predicción teórica: v2 escala casi linealmente mientras el cómputo manda y se estanca cuando la
comunicación lo alcanza.

\subsubsection{Número óptimo de procesos para v2}

Minimizando el modelo ajustado, $\dfrac{dT}{dp} = -\dfrac{a_1}{p^2} + \dfrac{a_2}{p\ln 2} + a_3 = 0$,
se obtiene la ecuación cuadrática $a_3\,p^2 + \dfrac{a_2}{\ln 2}\,p - a_1 = 0$, cuya raíz positiva es:
\[
p^* \approx 15.99 \;\approx\; 16.
\]
Este óptimo teórico coincide exactamente con el óptimo empírico $\arg\min_p T_p^{exp} = 16$
($T_{16} = 1.944$ s), lo que confirma que el modelo normalizado no solo reproduce la forma de la
curva sino también el punto de operación óptimo. Con el criterio alternativo de eficiencia
($E_p \geq 0.5$) el óptimo sería menor, pues para v2 la eficiencia cae por debajo de $0.5$ ya en
$p=16$.
